\documentclass[a4paper,10pt]{article}
\usepackage[utf8]{inputenc}
\usepackage[T1]{fontenc}
\usepackage[english]{babel}
\usepackage{enumitem}
\usepackage{hyperref}
\usepackage{geometry}
\geometry{margin=1in}

\hypersetup{
  pdftitle={Resume - Gian Pedro Rodrigues},
  pdfauthor={Gian Pedro Rodrigues},
  pdfsubject={Professional Resume},
  colorlinks=true,
  urlcolor=blue
}

\setlength{\parindent}{0pt}
\setlength{\parskip}{6pt}

\begin{document}

\vspace*{-1cm}
\begin{center}
    {\LARGE \textbf{Gian Pedro Rodrigues}}\\
    \vspace{0.2cm}
    \href{mailto:gianpedrodev@gmail.com}{gianpedrodev@gmail.com} \quad | \quad +55 (12) 99618-3274\\
    \href{https://github.com/GianSE}{github.com/GianSE} \quad | \quad 
    \href{https://linkedin.com/in/gian-pedro-rodrigues-330397163}{linkedin.com/in/gian-pedro-rodrigues}
\end{center}

\vspace{0.4cm}

\section*{Professional Summary}
Computer Engineering student at UTFPR with hands-on experience in Full Stack development (Node.js, React, Next.js) and projects involving Data and AI. Engaged in real-world applications through a junior enterprise, academic initiatives, and extracurricular activities.

\section*{Career Objective}
To work as a Full Stack Developer or Data Analyst, applying modern technologies to develop intelligent applications focused on delivering tangible results.

\section*{Education}
\begin{itemize}[leftmargin=*]
    \item \textbf{UTFPR – Federal University of Technology – Paraná} \hfill \textit{Jul 2022 – Jun 2027 (in progress)} \\
    Bachelor's Degree in Computer Engineering
    
    \item \textbf{ETEC – Technical School of the State of São Paulo} \hfill \textit{Jan 2015 – Dec 2017} \\
    Technical Degree in Accounting
\end{itemize}

\section*{Experiences}
\begin{itemize}[leftmargin=*]
    \item \textbf{Treasurer – UTFPR Student Council (DA)} \hfill \textit{2023 – 2024} \\
    Automated treasury processes using Google Sheets, App Script, and the PagBank API (payment gateway). Reduced manual spreadsheet work and improved financial control for the student organization.

    \item \textbf{Developer – LtSat UTFPR Project} \hfill \textit{2024 – 2025} \\
    Developed a data reception and visualization system for a satellite (STM32), with a back-end built in Node.js and interactive visualizations using D3.js on the front-end. Represented Brazil in the CanSat Competition (USA) as the country's only team.

    \item \textbf{Back-End Developer – Unect Junior Enterprise} \hfill \textit{2025 – present} \\
    Working with Node.js, Express.js, JWT, React.js, Next.js, and databases (MongoDB, PostgreSQL). The junior enterprise develops Web, Desktop, and Mobile application solutions.
\end{itemize}

\section*{Skills}
\begin{itemize}[leftmargin=*]
    \item \textbf{Back-End:} Node.js, Express.js, JWT, Java Spring Boot, MongoDB, PostgreSQL, MySQL
    \item \textbf{Front-End:} React.js, Next.js, HTML, CSS, Tailwind, Bootstrap
    \item \textbf{AI / Data:} Python, Pandas, NumPy, TensorFlow, Keras, Scikit-learn
    \item \textbf{Tools:} Git, GitHub, Linux, WSL, Postman, QtCreator
    \item \textbf{Other Languages:} C, C++
\end{itemize}

\section*{Languages}
\begin{itemize}[leftmargin=*]
    \item \textbf{English:} Intermediate (B1) – reading, writing, listening, and speaking. Studied at CALEM (Center for Modern Foreign Languages) – UTFPR.
\end{itemize}

\section*{Academic Projects}
\begin{itemize}[leftmargin=*]
    \item \textbf{Veterinary Clinic API} \\
    Developed a complete RESTful API using Node.js and Express, including JWT authentication and MongoDB Atlas integration. Used Postman for automated testing and organized request flows with scripts.\\  
    \href{https://github.com/GianSE/capacitacao-parte2}{github.com/GianSE/capacitacao-parte2}

    \item \textbf{UTF em Foco} \\
    Built a responsive informational website as an academic project using HTML, CSS (Bootstrap), and JavaScript. Modular structure and responsive layout for multiple devices.\\  
    \href{https://gianse.github.io/UTF_em_Foco/}{gianse.github.io/UTF_em_Foco/}

    \item \textbf{ClimaTech (Desktop)} \\
    Developed a desktop application in C++ (Qt Creator) integrating the OpenWeather API and PostgreSQL database to persist queried data.\\
    \href{https://github.com/GianSE/poo2_climatech}{github.com/GianSE/poo2_climatech}

    \item \textbf{MobileNet – Cifar10} \\
    Created a convolutional neural network model using MobileNet for image classification with the CIFAR-10 dataset. The project involved TensorFlow, Keras, preprocessing techniques, and performance evaluation in computer vision.\\    
    \href{https://colab.research.google.com/drive/1AGP831WCSEpkeN7zi-tSyno52o2E8LfE}{Google Colab Notebook}

    \item \textbf{Quantum Tunneling Simulation (Scientific Research)} \\
    Simulated the quantum tunneling phenomenon using Python. Modeled the Schrödinger equation for one-dimensional potentials, with graphical visualizations using Matplotlib and animations generated via FFmpeg. A project aimed at visually exploring quantum mechanics.\\  
    \href{https://colab.research.google.com/drive/122zQ2S9tb0PJ_tqczK8SwqvF77C_w4Md?usp=sharing}{Google Colab Notebook}
\end{itemize}

\end{document}
